\title{xLSTM: Extended Long Short-Term Memory}

\begin{abstract}
The foundational principles of the Long Short-Term Memory (LSTM), namely the constant error carousel and gating mechanisms, emerged in the 1990s and have since established LSTMs as resilient contributors to the evolution of deep learning, notably powering the earliest Large Language Models (LLMs). Nevertheless, the introduction of Transformer architectures, characterized by highly parallelizable self-attention, signaled a paradigm shift, allowing them to surpass LSTMs when deployed at large scale. In light of this development, we pose a straightforward question: To what extent can LSTMs be pushed in the realm of language modeling if they are scaled up to billions of parameters, drawing from innovations present in contemporary LLM frameworks while addressing longstanding LSTM drawbacks? Our investigation proceeds in two major directions. Firstly, we propose exponential gating enhanced with suitable normalization and stabilization methodologies. Secondly, we adapt the LSTM memory by devising: (i) sLSTM, which incorporates scalar memory, scalar updates, and novel methods for memory mixing; and (ii) mLSTM, distinguished by a fully parallelizable design through a matrix memory along with a covariance-based update process. We assemble these advances into xLSTM blocks using residual backbones, enabling the construction of deep xLSTM architectures via residual stacking. Collectively, exponential gating and altered memory architectures substantially elevate xLSTM models, enabling them to compare favorably—in both effectiveness and scalability—to state-of-the-art Transformer and State Space Models.\\
Code available at: \url{https://github.com/NX-AI/xlstm}
\end{abstract}

\section{Introduction}
\label{sec:intro}

The Long Short-Term Memory (LSTM) ideas~\citep{Hochreiter:91,Hochreiter:97nips,Hochreiter:97}, i.e., the constant error carousel and gating, were introduced to overcome the vanishing gradient problem of recurrent neural networks~\citep{Hochreiter:91,Hochreiter:00book}:

\begin{align}
\label{eq:lstm_idea}
  c_t \ = \ f_t \ c_{t-1} \ + \ 
          i_t \  z_t \ , \quad  
          h_t \ = \ o_t \ \psi({c_t}) \ . 
\end{align}

The cell state $c_{t-1}$ (shown in green) is adjusted by the additive integration of cell inputs $z_t$, with this process regulated through sigmoid-based gates (depicted in blue). The input gate $i_t$ and forget gate $f_t$ govern the manner and extent of this update, whereas the output gate $o_t$ determines the resulting output from the memory cell, specifically the hidden state $h_t$. Before producing the hidden state, the cell state undergoes normalization or compression via the function $\psi$, followed by gating at the output stage.

LSTMs have achieved significant success across multiple application areas \citep{Hochreiter:01,Hochreiter:07,Schmidhuber:15}, dominating the landscape of text generation prior to the advent of Transformers in 2017~\citep{Vaswani:17}. Their proficiency has been empirically validated in a broad array of sequence-dependent tasks, including but not limited to, text generation~\citep{Graves:13,Karpathy:15blog}, handwriting synthesis~\citep{Graves:13}, machine translation using sequence-to-sequence frameworks~\citep{Sutskever:14nips}, program evaluation~\citep{Zaremba:14arxiv}, image caption creation~\citep{Karpathy:15,Hossain:19}, automated code synthesis~\citep{Karpathy:15blog}, rainfall-runoff prediction~\citep{Kratzert:18,Kratzert:19}, and hydrological forecasting for flood warnings~\citep{Nearing:24}. Within the domain of reinforcement learning, LSTMs have established themselves as leading sequence modeling architectures, powering systems such as AlphaStar for StarCraft II~\citep{Vinyals:17short}, OpenAI Five for Dota 2~\citep{OpenAI:19}, and magnetic plasma controllers for nuclear fusion~\citep{Degrave:22short}. LSTMs are notably effective at learning high-level abstractions, skillfully identifying and retaining semantic features within their memory mechanisms~\citep{Karpathy:15blog}, a property illustrated by the discovery of units sensitive to numeric and syntactic patterns~\citep{Lakretz:19}, linguistic properties~\citep{Bau:19}, and sentiment~\citep{Radford:17arxiv}. Their enduring relevance in practical applications~\citep{Degrave:22short,Nearing:24} demonstrates their long-term robustness and utility.

Although LSTMs have achieved remarkable results, they exhibit three significant shortcomings: (i) Inability to modify storage choices retrospectively. This constraint is illustrated using the {\it Nearest Neighbor Search} task (further details are provided in Appendix~\ref{sec:appNearestNeighborSearch}). Here, one must examine a sequence to find the vector closest to a reference and return its associated value at the end of the sequence. The left-hand side of Figure~\ref{fig:lstmProblems} presents the mean squared error for this scenario. LSTM networks are notably deficient at updating stored values once a more similar vector appears, a shortcoming that the proposed xLSTM addresses through exponential gating mechanisms. (ii) Constrained storage capability, necessitating the condensation of information into singular scalar cell states. This is demonstrated by the {\it Rare Token Prediction} task, depicted in the right-hand panel of Figure~\ref{fig:lstmProblems}, where perplexity scores for token prediction on the Wikitext-103~\citep{Merity:17} corpus are shown across frequency bins. LSTMs underperform on infrequent tokens due to these restricted memory resources, while the xLSTM leverages a matrix-based memory to overcome this obstacle. (iii) Absence of parallelism due to memory interdependency, stemming from recurrent connections between hidden states over time, which impose strictly sequential computation. 

The emergence of Transformers~\citep{Vaswani:17} in language modeling has, in part, been driven by these intrinsic limitations of the LSTM architecture. This raises the question: by eliminating these constraints and scaling LSTMs to sizes comparable to present-day Large Language Models, what new levels of performance in language modeling might be realized?

\section{Extended Long Short-Term Memory}
\label{sec:xLSTM}

In response to the constraints inherent in standard LSTMs, Extended Long Short-Term Memory (xLSTM) incorporates two significant advancements beyond the framework outlined in Equation~\eqref{eq:lstm_idea}. These enhancements—namely exponential gating and the development of novel memory structures—expand the family of LSTM models with two additions: (i) sLSTM, detailed in Section~\ref{sec:sLSTM}, which utilizes a scalar-valued memory, performs scalar updates, and employs a memory mixing mechanism; and (ii) mLSTM, presented in Section~\ref{sec:mLSTM}, which leverages a matrix-based memory in conjunction with a covariance-driven (outer product) update procedure that enables full parallelism. Both sLSTM and mLSTM build on LSTM by employing exponential gating mechanisms. Notably, in pursuit of parallelizable operations, mLSTM eliminates memory mixing, specifically the recurrent connections among hidden states. Both mLSTM and sLSTM architectures can be generalized to encompass multiple memory cells, where sLSTM also supports memory mixing between cells. Moreover, sLSTM accommodates a multi-head architecture, wherein memory mixing is absent across distinct heads but permitted among cells within each head. This design—introducing head-based modularity in sLSTM combined with exponential gating—offers a new paradigm for organizing memory mixing. For mLSTM, configurations with multiple heads and multiple cells are functionally interchangeable.

These newly conceived LSTM variants can be embedded within residual block modules to form xLSTM blocks (refer to Section~\ref{sec:xLSTMblock}). By arranging such xLSTM blocks residually within network architectures, comprehensive xLSTM architectures can be constructed (see Section~\ref{sec:xLSTMarchitecure}). Figure~\ref{fig:xlstm_sketch} provides a schematic depiction of the xLSTM framework and its constituent elements.

\subsection{Review of the Long Short-Term Memory}
\label{sec:orgLSTM}

The foundational concept behind LSTM~\citep{Hochreiter:91,Hochreiter:97nips,Hochreiter:97} established a scalar memory cell, serving as a core mechanism for computation and storage. This architecture mitigates the problem of vanishing gradients~\citep{Hochreiter:91,Hochreiter:00book} via the constant error carousel, which preserves information during cell state updates. The memory cell operates with three distinct gates: the input, output, and forget gates, with the forget gate being proposed by \citet{Gers:00}. The LSTM memory cell is updated at time step $t$ according to the following rules:

\begin{align}
c_t \ &= \  f_t \ c_{t-1} \ + \ i_t \ z_t &  & &\text{cell state} \\
h_t  \ &= \ o_t \ \tilde{h}_t \ , & \tilde{h}_t \ &= \ \psi \left( c_t\right)
  &\text{hidden state} \\
z_t \ &= \ \varphi \left( \tilde{z}_t \right) \ , 
  &\tilde{z}_t \ &=  \ \mathbf{w}^\top_{z} \ \mathbf{x}_t \ + \
  r_{z}  h_{t-1} \ + \  b_{z} \ \
  &\text{cell input} \\
i_t \ &= \ \sigma \left( \tilde{i}_t  \right) \ , 
  &\tilde{i}_t \ &= \ \mathbf{w}^\top_{i} \ \mathbf{x}_t \ + \
  r_{i}  \ h_{t-1} \ + \  b_{i} \ \
  &\text{input gate} \\
f_t \ &= \ \sigma \left(  \tilde{f}_t \right) \ , 
  &\tilde{f}_t \ &= \ \mathbf{w}^\top_{f} \ \mathbf{x}_t  \ + \
  r_{f}  \ h_{t-1} \ + \  b_{f} \ \
  &\text{forget gate} \\
o_t \ &= \ \sigma \left( \tilde{o}_t \right) \ , 
  &\tilde{o}_t  \ &= \ \mathbf{w}^\top_{o} \ \mathbf{x}_t \ + \
  r_{o}  \ h_{t-1} \ + \  b_{o} \ \
  &\text{output gate} 
\end{align}

The input weight vectors $\mathbf{w}_{z}$, $\mathbf{w}_{i}$, $\mathbf{w}_{f}$, and $\mathbf{w}_{o}$ link the inputs $\mathbf{x}_t$ to the cell input, input gate, forget gate, and output gate, respectively. Similarly, the recurrent weights $r_{z}$, $r_{i}$, $r_{f}$, and $r_{o}$ connect the hidden state $h_{t-1}$ to the corresponding components: cell input, input gate, forget gate, and output gate. The parameters $b_{z}$, $b_{i}$, $b_{f}$, and $b_{o}$ serve as biases for these respective connections. The activation functions for the cell input and hidden state are denoted by $\varphi$ and $\psi$ (commonly chosen as $\tanh$). The function $\psi$ is particularly important for compressing the values of the cell state, thus preventing it from growing without bound. Each of the gates utilizes a sigmoid activation, specifically,
$\sigma \left( x \right)= 1/(1+\exp(-x))$.
Subsequent advances grouped multiple scalar memory cells, $\mathbf{c}_t \in \mathbb{R}^{d}$, into a vector form. This reformulation enables the use of recurrent weight matrices $\mathbf{R} \in \mathbb{R}^{d \times d}$, allowing interactions among the outputs of different memory cells~\citep{Greff:15}; further exposition can be found in Appendix~\ref{sec:apporgLSTM}. Findings from ablation experiments confirmed that each part of the memory cell plays an indispensable role~\citep{Greff:15}.

\subsection{sLSTM}
\label{sec:sLSTM}

In order to enhance LSTMs with flexible memory updating, we implement exponential gates (highlighted in red), incorporating both normalization and stabilization mechanisms. Specifically, input and forget gates are designed to utilize exponential activation functions. As part of the normalization process, we add a normalizer state that accumulates the product of the input gate and all subsequent forget gates.

The scalar sLSTM forward pass is:

\begin{align}
\label{eq:slstmforward}
c_t \ &= \  f_t \ c_{t-1} \ + \ i_t \ z_t & & 
 &\text{cell state} \\
n_t \ &= \  f_t \ n_{t-1} \ + \ i_t  & & 
 &\text{normalizer state} \\
h_t \ &= \ o_t \ \tilde{h}_t \ ,
  &\tilde{h}_t \ &= \ c_t / n_t
&\text{hidden state} \\
z_t \ &= \ \varphi \left( \tilde{z}_t \right) \ , 
  &\tilde{z}_t \ &=  \ \mathbf{w}^\top_{z} \ \mathbf{x}_t \ + \
  r_{z}  h_{t-1} \ + \  b_{z}
  &\text{cell input} \\
i_t \ &= \ \!\exp \left( \tilde{i}_t  \right) \ , 
  &\tilde{i}_t \ &= \ \mathbf{w}^\top_{i} \ \mathbf{x}_t \ + \
  r_{i}  \ h_{t-1} \ + \  b_{i}
  &\text{input gate} \\
f_t \ &= \ \sigma \left(  \tilde{f}_t \right) \ \text{OR} \
\!\exp \left(  \tilde{f}_t \right) \ ,
  &\tilde{f}_t \ &= \ \mathbf{w}^\top_{f} \ \mathbf{x}_t  \ + \
  r_{f}  \ h_{t-1} \ + \  b_{f}
  &\text{forget gate} \\
o_t \ &= \ \sigma \left( \tilde{o}_t \right) \ , 
  &\tilde{o}_t  \ &= \ \mathbf{w}^\top_{o} \ \mathbf{x}_t \ + \
  r_{o}  \ h_{t-1} \ + \  b_{o}
  &\text{output gate} 
\end{align}

We transfer the original LSTM gating techniques, i.e., input- and/or hidden-dependent gating plus bias term, to the new architectures. Exponential activation functions can lead to large values that cause overflows. Therefore, we stabilize gates with an additional state \mbox{$m_t$~\citep{Milakov:18arxiv}}:

\begin{align}
\label{eq:slstmstabil}
m_t \ &= \ \max \left( \log ( f_t ) + m_{t-1} , \log ( i_t ) \right) &\text{stabilizer state} \\
i'_t \ &= \ \!\exp \left( \log \left ( i_t \right) - m_t \right) \ = \ \!\exp \left( \tilde{i}_t - m_t \right) \ \, 
  &\text{stabil. input gate} \\
f'_t \ &= \ \!\exp \left( \log \left( f_t \right) + m_{t-1} - m_t \right) \ \, 
  &\text{stabil. forget gate}
\end{align}

As demonstrated in Appendix~\ref{sec:appsLSTM}, substituting $f_t$ with $f'_t$ and $i_t$ with $i'_t$ in the forward propagation does not alter either the final network output or the gradients of the loss function with respect to the model parameters.

\paragraph{Emergent Memory Mixing.} Similar to the original LSTM, sLSTM accommodates several memory cells, as detailed in Appendix~\ref{sec:appsLSTM}. The presence of multiple memory cells allows for memory mixing that is facilitated by recurrent pathways $\mathbf{R}_{\mathbf{z}}$, $\mathbf{R}_{\mathbf{i}}$, $\mathbf{R}_{\mathbf{f}}$, $\mathbf{R}_{\mathbf{o}}$ linking the hidden state vector $\mathbf{h}$ to the memory cell input $\mathbf{z}$ and gate variables $\mathbf{i}$, $\mathbf{f}$, $\mathbf{o}$. A distinctive characteristic of memory mixing in this context is the involvement of exponential gating. The novel sLSTM architecture supports several heads, permitting memory mixing only within individual heads, not between them. Integrating exponential gating and head structures in sLSTM introduces a fundamentally new paradigm for memory mixing.

\subsection{mLSTM}
\label{sec:mLSTM}

To improve the memory capabilities of LSTMs, we generalize the LSTM cell state from a single value $c \in \mathbb{R}$ to a matrix $\mathbf{C} \in \mathbb{R}^{d \times d}$. As a result, retrieval is implemented using matrix multiplication. At each time step $t$, the model records a vector pair: the key $\mathbf{k}_t \in \mathbb{R}^d$ and the value $\mathbf{v}_t \in \mathbb{R}^d$, following the nomenclature of the Transformer architecture. Subsequently, at a later time $t + \tau$, the value $\mathbf{v}_t$ is recovered with a query vector $\mathbf{q}_{t+\tau} \in \mathbb{R}^d$. This framework is equivalent to the Bidirectional Associative Memories (BAMs) paradigm \citep{Kohonen:72,Anderson:72,Nakano:72,Anderson:77}.

The covariance update rule~\citep{Sejnowski:77,Dayan:91} for storing a key-value pair is

\begin{align}
\mathbf{C}_t \ &= \ \mathbf{C}_{t-1} \ + \ \mathbf{v}_{t} \ \mathbf{k}_t^\top \ .
\end{align}

We postulate that a layer-norm is applied prior to mapping inputs to keys and values, which ensures they possess zero mean. The covariance update mechanism achieves optimality~\citep{Dayan:91} by maximizing the discriminability among retrieved binary vectors; this objective aligns with optimizing the signal-to-noise ratio. Enhanced separability can be realized if retrieval is confined to pairwise interactions, accepting the resultant quadratic time complexity characteristic of attention~\citep{Krotov:16,Krotov:17,Ramsauer:21}. The covariance update strategy mirrors that of Fast Weight Programmers \citep{Schmidhuber:92ncfastweights,Schlag:21}, which have subsequently introduced a fixed decay coefficient for $\mathbf{C}_{t-1}$ and a fixed learning rate applied to $\mathbf{v}_{t} \mathbf{k}_t ^\top$~\citep{Ba:16}. Following this rationale, we embed the covariance update process within the LSTM architecture: the forget gate modulates the decay, the input gate governs the learning rate, and the output gate controls the amplitude of the retrieved vector.

For this matrix memory, the normalizer state is the weighted sum of key vectors, where each key vector is weighted by the input gate and all future forget gates. Again, the normalizer state keeps record of the strength of the gates. Since the dot product between query and normalizer state can be close to zero, we use the absolute value of this dot product and lower bound it by a threshold (typically 1.0) as done previously~\citep{Sun:23arxiv}. The mLSTM forward pass is:

\begin{align}
\label{eq:mlstm_recurrent_begin}
\mathbf{C}_t \ &= \  f_t \ \mathbf{C}_{t-1} \ + \ 
  i_t \ \mathbf{v}_t \ \mathbf{k}_t^\top & &  &\text{cell state} \\
\mathbf{n}_t \ &= \  f_t \ \mathbf{n}_{t-1} \ + \ 
  i_t \ \mathbf{k}_t & &  &\text{normalizer state} \\
\mathbf{h}_t  \ &= \ \mathbf{o}_t \ \odot \ \tilde{\mathbf{h}}_t
\ , \qquad \qquad 
\tilde{\mathbf{h}}_t \ = \   \mathbf{C}_t \mathbf{q}_t \ / \ 
  \max \left\{ |\mathbf{n}_t^\top \mathbf{q}_t|, 1 \right\} 
   &&&\text{hidden state} \\
\mathbf{q}_t \ &= \ \mathbf{W}_q \ \mathbf{x}_t \ + \ \mathbf{b}_q  & &  &\text{query input} \\
\mathbf{k}_t \ &= \ \frac{1}{\sqrt{d}} \mathbf{W}_k \ \mathbf{x}_t \ + \ \mathbf{b}_k  & &  &\text{key input} \\
\mathbf{v}_t \ &= \ \mathbf{W}_v \ \mathbf{x}_t \ + \ \mathbf{b}_v  & &  &\text{value input} \\
i_t \ &= \ \!\exp \!  \! \left( \tilde{i}_t  \right) \ , 
 \qquad \qquad \ \ \ \ \,
  \tilde{i}_t \ = \ \mathbf{w}^\top_{i} \ \mathbf{x}_t \ + \  b_{i} \ \ \ 
  &&&\text{input gate} \\
f_t \ &= \ \sigma \!  \left(  \tilde{f}_t \right) \ \text{OR} \ \!\exp \!  \! \left(  \tilde{f}_t \right) , \ \ \: \!
\tilde{f}_t \ = \ \mathbf{w}^\top_{f} \ \mathbf{x}_t  \ + \
  b_{f}
  &&& \text{forget gate} \\
\label{eq:mlstm_recurrent_end}
\mathbf{o}_t \ &= \ \sigma \left( \tilde{\mathbf{o}}_t \right) \ , \qquad \qquad \quad \ \ \ \,
  \tilde{\mathbf{o}}_t  \ = \ \mathbf{W}_{\mathbf{o}} \ \mathbf{x}_t \ + \
  \mathbf{b}_{\mathbf{o}}
  &&&\text{output gate} 
\end{align}

The mLSTM architecture allows for several memory cells, similar to the classic LSTM model. In mLSTM, the distinction between multiple heads and multiple cells becomes negligible due to the absence of memory mixing mechanisms. To ensure stability in the exponential gates of mLSTM, we apply the stabilization methods utilized for sLSTM (refer to Equation~\ref{eq:slstmstabil}). Given the lack of memory mixing in mLSTM, the recurrent process can be expressed in a parallel form. Further technical details are provided in Appendix~\ref{sec:appmLSTM}.

\subsection{xLSTM Architecture}
\label{sec:xLSTMarch}

\paragraph{xLSTM Blocks.}
\label{sec:xLSTMblock}
An xLSTM block aims to non-linearly encode prior sequence information into a high-dimensional representation, which enhances the model's ability to distinguish among various contexts or historical sequences. Such differentiation of histories is essential for accurate prediction of subsequent elements (e.g., the next token) in a sequence. In light of Cover’s Theorem~\citep{Cover:65}, which asserts that patterns embedded non-linearly into higher-dimensional spaces are more likely to be linearly separable than those in the original input space, our approach is motivated by this principle. We investigate two types of residual block structures: (i) a residual block utilizing post up-projection (akin to Transformer models), where historical information is first non-linearly integrated in the initial space, followed by a linear projection into a higher-dimensional space, application of a non-linear activation, and then a linear projection back into the original dimension; refer to the left illustration in Figure~\ref{fig:Backbones} and the third column of Figure~\ref{fig:xlstm_sketch}. A comprehensive depiction is available in Figure~\ref{fig:appsLSTM_detailed} in the appendix. (ii) a residual block employing pre up-projection (similar to State Space Models), which performs a linear mapping into a higher-dimensional space,

It performs a non-linear aggregation of preceding information within a high-dimensional space, followed by a linear transformation that projects back into the original space. In the configuration where the xLSTM block incorporates an sLSTM, we employ a post up-projection module. Conversely, for xLSTM blocks utilizing an mLSTM, we apply the pre up-projection module due to the increase in memory capacity inherent to the higher-dimensional representation. Further clarification is available in the left panel of Figure~\ref{fig:Backbones}, the third column of Figure~\ref{fig:xlstm_sketch}, as well as in Figure~\ref{fig:appsLSTM_detailed} in the appendix.

\paragraph{xLSTM Architecture. }
\label{sec:xLSTMarchitecure}
The xLSTM model is formed by stacking modular components in a residual manner, as outlined by previous work~\citep{Srivastava:15,He:16}. For this, we adopt the prevalent pre-LayerNorm~\citep{Ba:16b} residual backbone architecture, which is standard in modern Large Language Models. An illustration is provided in the final column of Figure~\ref{fig:xlstm_sketch}.

\subsection{Memory and Speed Considerations}

In contrast to Transformers, xLSTM networks maintain linear computational cost and fixed memory usage relative to sequence length. Due to their compressive memory properties, xLSTM models are particularly advantageous for industrial deployment and edge device implementation.

The memory mechanism in mLSTM operates without requiring additional parameters, but the associated computational burden stems from its $d \times d$ memory matrix and corresponding $d \times d$ updates. This represents a balance between memory size and computational workload. Nonetheless, these operations can be efficiently parallelized on GPUs, making their contribution to overall wall-clock time relatively minor.

Although mLSTM lends itself to parallel processing in a manner similar to approaches such as FlashAttention~\citep{Dao:22,Dao:23} and GLA~\citep{Yang:23arxiv}, sLSTM lacks parallelizability because of the involvement of memory mixing through hidden-hidden interactions. To address this, we implemented an efficient CUDA-based approach with register-level GPU memory optimizations; as a result, its runtime is typically less than twice as long as that of mLSTM.

\section{Related Work}
\label{sec:related}

\paragraph{Linear Attention.} Numerous strategies have been proposed to address the challenge of the Transformer’s quadratic computational complexity relative to context length, aiming to achieve linear attention scaling. The Synthesizer generates attention scores synthetically, eliminating explicit token-to-token dependencies~\citep{Tay:20arxiv}. Linformer implements self-attention using a low-rank matrix, thereby facilitating linear approximation~\citep{Wang:20arxiv}. Linear Transformer reformulates the attention operation in a linear manner \citep{Katharopoulos:20}. Performer introduces a linear approximation to the attention softmax function using a method based on positive orthogonal random features~\citep{Choromanski:21}. Additionally, the Structured Global Convolution (SGConv)~\citep{Li:22} and the Hyena Hierarchy~\citep{Poli:23} replace attention with rapid, extensive convolutions.

\paragraph{State Space Models.} In recent developments, State Space Models (SSMs) have garnered significant attention due to their linear complexity with respect to context length and competitive performance relative to Transformers. Early advancements in this area include the introduction of the Structured State Space sequence model (S4)~\citep{Gu:21}, followed by notable successors such as the Diagonal State Space (DSS) model~\citep{Gupta:22}, Gated State Space (GSS) models~\citep{Mehta:22}, S5 model~\citep{Smith:22}, Bidirectional Gated SSM (BiGS)~\citep{Wang:22}, H3 model~\citep{Fu:23}, and Mamba~\citep{Gu:24arxiv}.

\paragraph{Recurrent Neural Networks.} In efforts to overcome limitations of Transformers and attention mechanisms regarding sequence length, Recurrent Neural Networks (RNNs) have emerged as viable alternatives thanks to their linear scaling with context size. Deep Linear Recurrent Units (LRUs) within RNN architectures have produced strong results in language modeling tasks~\citep{Orvieto:23, De:24arxiv}. Similarly, innovations like Hierarchically Gated Linear RNN (HGRN)~\citep{Qin:23} and its subsequent iteration, HGRN2~\citep{Qin:24arxiv}, have demonstrated efficacy in this domain. Among prominent RNN-based strategies for large language models is RWKV~\citep{Peng:23arxivshort, Peng:24arxivshort}, which has achieved performances on par with Transformer models.

\paragraph{Gating.}
LSTM’s foundational concept of gating has not only persisted but also found renewed application and varied interpretation in contemporary research. Numerous recent models leverage gating mechanisms, including HGRN~\citep{Qin:23}, HGRN2~\citep{Qin:24arxiv}, Gated Linear Attention (GLA)~\citep{Yang:23arxiv}, Gated State Space (GSS) models~\citep{Mehta:22}, Bidirectional Gated SSM (BiGS)~\citep{Wang:22}, Moving Average Equipped Gated Attention (MEGA)~\citep{Ma:22}, RWKV~\citep{Peng:23arxivshort}, and Mamba~\citep{Gu:24arxiv}.

\paragraph{Covariance Update Rule.} 
To bolster the memory capacity, we integrated a covariance update rule within the matrix memory of the mLSTM cell. Similar update procedures are present in other methods such as Fast Weight Programmers \citep{Schmidhuber:92ncfastweights,Schlag:21}, RWKV-5 and RWKV-6~\citep{Peng:24arxivshort}, Retention~\citep{Sun:23arxiv}, Linear Transformer~\citep{Katharopoulos:20}, and HGRN2~\citep{Qin:24arxiv}.

\paragraph{Most Related.} The models most analogous to xLSTM in their conceptual foundations are Retention~\citep{Sun:23arxiv}, RWKV~\citep{Peng:23arxivshort,Peng:24arxivshort}, and HGRN2~\citep{Qin:24arxiv}, which implement matrix memory and gating mechanisms. Nonetheless, unlike the proposed sLSTM, these methods do not incorporate memory mixing. The presence of memory mixing is crucial for addressing state tracking challenges, positioning LSTMs as more capable frameworks compared to State Space Models (SSMs) and Transformers~\citep{Merrill:24,Deletang:23}. Effective state tracking is necessary for tasks such as code evaluation and entity monitoring within extensive narratives.

\paragraph{Residually Stacking Architectures.} The xLSTM models, in alignment with most state-of-the-art deep learning systems, employ a residually stacked arrangement of modules~\citep{Srivastava:15,He:16}.

This methodology paved the way for very deep convolutional architectures~\citep{He:16} and was central to the development of Transformer models~\citep{Vaswani:17}. Transformers serve as the foundational paradigm underlying Large Language Models (LLMs) such as GPT-3~\citep{Brown:20short}, ChatGPT~\citep{Schulman:22short}, GPT-4~\citep{Achiam:23gpt4short}, Megatron-LM~\citep{Shoeybi:19}, Gopher~\citep{Rae:21short}, ERNIE 3.0 Titan~\citep{Wang:21arxivshort}, GLaM~\citep{Du:21glamshort}, Chinese M6~\citep{Lin:21}, AlexaTM 20B (multilingual)~\citep{Soltan:22}, OPT~\citep{Zhang:22opt}, Chinchilla~\citep{Hoffmann:22short}, BLOOM~\citep{Scao:22arxivshort}, GLM-130B~\citep{Zeng:22arxivshort}, LaMDA~\citep{Thoppilan:22short}, PaLM~\citep{Chowdhery:22short}, Llama~\citep{Touvron:23arxiv}, and Gemini~\citep{GeminiTeam:23arxiv,Reid:24arxiv}.

\section{Experiments}
\label{sec:Exp}

This section presents an empirical analysis of xLSTM, emphasizing its performance relative to established approaches in the domain of language modeling. Section~\ref{sec:ExpSyntethic} explores xLSTM's strengths and weaknesses through synthetic task evaluations. In Section~\ref{sec:ExpComparison}, we measure and contrast the validation perplexities achieved by several contemporary language modeling models, each trained on a dataset containing 15B tokens from SlimPajama~\citep{Soboleva:23}. Additionally, we perform ablation experiments on xLSTM using this dataset. The scaling properties of these models are subsequently examined, following methodologies similar to those described in \citet{Kaplan:20} and \citet{Brown:20short}.

Section~\ref{sec:ExpLanguage} advances to a comprehensive language modeling study. Here, xLSTM is compared with the top-performing baselines identified in Section~\ref{sec:ExpComparison}, after training on an extended corpus of 300B tokens from SlimPajama~\citep{Soboleva:23}. We analyze model performance in several dimensions: first, the ability to generalize to longer contextual sequences; second, using metrics of validation perplexity and outcomes in downstream tasks~\citep{Sutawika:24}; third, evaluating across the 571 distinct domains provided by the PALOMA language benchmark~\citep{Magnusson:23arxivshort}; and fourth, reassessing scaling dynamics, this time employing a training set that is twenty times larger.

\label{sec:xlstm_blocks_notation}
Throughout our experiments, we denote xLSTM[$a$:$b$] as the configuration where the fraction $a/b$ represents the number of mLSTM-based to sLSTM-based xLSTM blocks. For instance, xLSTM[7:1] corresponds to a setup with eight blocks, of which seven utilize mLSTM and one employs sLSTM. With a typical configuration involving 48 blocks in total, this yields 42 mLSTM-based blocks and 6 sLSTM-based blocks. Additionally, every experimental setting employs up-projection blocks before mLSTM modules and after sLSTM modules, respectively.

\subsection{Synthetic Tasks and Long Range Arena}
\label{sec:ExpSyntethic}
Initially, we evaluate how xLSTM’s novel exponential gating combined with memory mixing performs on formal languages~\citep{Deletang:23}. Subsequently, we examine xLSTM’s new matrix memory capabilities through the Multi-Query Associative Recall task~\citep{Arora:23arxiv}. Lastly, we analyze xLSTM’s handling of extended sequence data by benchmarking it on the Long Range Arena~\citep{Tay:21}.

\paragraph{Evaluation of xLSTM's Exponential Gating with Memory Mixing.} We evaluate the effectiveness of xLSTM's novel exponential gating mechanism augmented with memory mixing, which is hypothesized to facilitate the resolution of state tracking tasks~\citep{Merrill:24, Merrill:23}. Building on the formal language benchmarks introduced by \citet{Deletang:23}, we adapt these tasks to support training across variable sequence lengths to allow for extrapolation to previously unseen lengths. A comprehensive explanation of each task, along with more extensive experimental outcomes, is included in Appendix~\ref{sec:appExpSynthetic-formalLang}. For benchmarking, we situate xLSTM among competing frameworks such as Transformers, State Space Models, and various Recurrent Neural Networks. Performance is assessed by measuring accuracy restricted to tokens pertinent to each respective task, with the scoring normalized between 0 (chance level) and 1 (exact performance). Our comparison considers 2-block versions of several architectures, specifically: xLSTM[0:1] (containing only sLSTM units), xLSTM[1:0] (exclusively mLSTM units), xLSTM[1:1], Llama, Mamba, RWKV, Retention, Hyena, LSTM, and LSTM implemented within Transformer blocks (LSTM (Block)). Experimental results are visualized in Figure~\ref{fig:formal-main}. It is observed that models such as Transformers and State Space Models, when lacking memory mixing (and thus proper state tracking), are unable to address certain regular grammars, such as in the parity task. This observation corroborates prior evidence indicating that Transformers and State Space Models possess fundamentally less computational capacity for state tracking than RNNs~\citep{Merrill:24, Merrill:23, Deletang:23}.

\paragraph{Test of xLSTM's Memory Capacities on Associative Recall Tasks.} This study evaluates the matrix memory mechanism of xLSTM by assessing its memory limits using the Multi-Query Associative Recall task~\citep{Arora:23arxiv}. Within each input sequence, key-value pairs are randomly selected from an extensive vocabulary and are required to be stored for subsequent retrieval. To create a more demanding version of the original task, we scale up both the number of key-value pairs, reaching as many as 256, and the overall context length, extending it to 2048. This configuration enables a more comprehensive examination of memory capacities across various models. Specifically, we benchmark 2-block variants of Llama, Mamba, RWKV-5, RWKV-6, xLSTM[1:1], and xLSTM[1:0], measuring performance by the retrieval accuracy of key-value associations. Transformers (like Llama), characterized by memory that scales exponentially with the representation dimension~\citep{Ramsauer:21}, serve as the reference point for optimal performance on this benchmark. Comparative outcomes are displayed in Figure~\ref{fig:mqar-main}.

Among all models that are not based on the Transformer architecture, xLSTM[1:1] delivers superior performance, even in smaller-scale configurations. Notably, the inclusion of the sLSTM block does not impede, but rather appears to facilitate, memory retention—particularly apparent in the most challenging scenario involving 256 pairs. Additional findings can be found in Appendix~\ref{sec:appExpSynthetic-mqar}, where extrapolation experiments reveal that the improved memory of xLSTM permits generalization to sequence contexts longer than those encountered during model training.

\paragraph{Test of xLSTM's Long Context Capabilities on Long Range Arena.} In evaluating xLSTM's effectiveness with extended sequences and substantial context lengths, various techniques are benchmarked using the Long Range Arena~\citep{Tay:21}. Across all tasks, xLSTM persistently achieves high results, indicating that the architecture is particularly adept at addressing diverse challenges associated with processing long-range context.

For more details, see Appendix~\ref{sec:appExpSynthetic-lra}.

\subsection{Method Comparison and Ablation Study}
\label{sec:ExpComparison}

In order to investigate our central research question—namely, the potential capabilities of the proposed LSTM variants when expanded for language modeling tasks—we conduct experiments by training xLSTMs, Transformers, State Space Models, along with additional approaches, on 15B tokens sourced from SlimPajama using a uniform auto-regressive framework. Following training, we evaluate all models on the validation dataset and undertake ablation analyses specifically for xLSTMs.

\paragraph{Performance Evaluation of xLSTM Against Alternative Approaches.} To facilitate comparison, we trained each model using 15B tokens derived from SlimPajama~\citep{Soboleva:23}. The assessment was conducted by measuring each model’s perplexity on the validation data. The models incorporated into this comparison include: the newly proposed xLSTM, GPT-3 (Transformer)~\citep{Brown:20short}, Llama (Transformer)~\citep{Touvron:23arxiv}, H3 (SSM)~\citep{Fu:23}, Mamba (SSM)~\citep{Gu:24arxiv}, RWKV-4 (RNN)~\citep{Peng:23arxivshort}, RWKV-5 (RNN)~\citep{Peng:24arxivshort}, RWKV-6 (RNN)~\citep{Peng:24arxivshort}, GLA (linear Transformer)~\citep{Yang:23arxiv}, HGRN (RNN)~\citep{Qin:23}, HGRN2 (RNN)~\citep{Qin:24arxiv}, RetNet (linear Transformer)~\citep{Sun:23arxiv}, Hyena (linear Transformer)~\citep{Poli:23}, xLSTM[1:0], and xLSTM[7:1] (see Section~\ref{sec:xlstm_blocks_notation}). 

Training employed mixed precision across the board; however, for RWKV-5, RWKV-6, GLA, and HGRN2, the mixed precision procedure was not managed through PyTorch’s built-in automatic mixed precision functionality (refer to Appendix Section~\ref{sec:appGenTrainProc} for further details). The comparison groups the models as follows: (a) Transformers, (b) State Space Models (SSMs), and (c) Recurrent Neural Networks (RNNs), along with linear Transformers—where linear Transformers adopt a linear method to replace traditional attention in the Transformer architecture. All models were configured to mirror the scale of a GPT-3 model with 350M parameters, meaning they used embedding size 1024 and 24 residual blocks. Notably, GPT-3 exclusively employs weight sharing between the input token embedding and the output embedding, resulting in a lower parameter count.

As seen in Table~\ref{tab:spaj15b_model_results}, xLSTM achieves the best performance, surpassing all previous methods with respect to validation perplexity. Additional specifics are provided in Appendix~\ref{sec:appExpComparison}. Figure~\ref{fig:temp_sclaw15B} further illustrates model scaling within this context, demonstrating that xLSTM’s advantage persists as model sizes increase.

\paragraph{Ablation Studies.}
As shown in Table~\ref{tab:spaj15b_model_results} and Figure~\ref{fig:temp_sclaw15B}, xLSTM delivers strong language modeling performance when trained with 15B tokens from SlimPajama. To systematically examine the impact of each modification from LSTM to xLSTM, we incrementally transform the standard LSTM architecture into the xLSTM design. The initial step involves incorporating LSTM layers within residual backbones equipped with pre-LayerNorm. Next, we introduce a post up-projection block to the model. The process concludes with the addition of exponential gating mechanisms and matrix memory components.

The results are shown in Table~\ref{tab:ablstudies} (top). 

The ablation experiments reveal that the notable gains in performance are a result of the combined impact of exponential gating and matrix memory. Given the pivotal role of gating within RNNs and State Space Models, we further investigate various gating strategies by conducting ablations. As shown in Table~\ref{tab:ablstudies} (bottom), our analysis demonstrates that allowing each gate to be learnable and directly affected by the input leads to consistent improvements. More detailed examinations of the underlying backbone components are presented in Appendix~\ref{sec:appAblDetails}.

\subsection{xLSTM as Large Language Model}
\label{sec:ExpLanguage}

In this work, we conduct comprehensive language modeling experiments at scale, exploring the viability of xLSTM as a large language model (LLM). To do so, we expand the training dataset to encompass 300B tokens sourced from SlimPajama. This training regime mirrors the token count utilized in Mamba~\citep{Gu:24arxiv} and Griffin~\citep{De:24arxiv}.

For our comparative analysis, we benchmark xLSTM against RWKV-4, Llama, and Mamba—each representing a distinct category as detailed in Section~\ref{sec:ExpComparison}. RWKV-4 is chosen to stand for the RNN class, since suitable training precision for RWKV-5, RWKV-6, and HGRN2 was established only after commencement of the 300B-token runs (refer to Appendix~\ref{sec:appGenTrainProc}).

We train models across several parameter scales (125M, 350M, 760M, 1.3B), examine length extrapolation properties, and evaluate their validation set performance. Additionally, we scrutinize their abilities in downstream tasks, measure their effectiveness in language modeling across the 571 text domains found in the PALOMA benchmark, and perform a scaling law analysis.

\paragraph{Sequence Length Extrapolation.}
\label{sec:spaj300B_ctx_extrapolation}
To begin, we investigate how well 1.3B-parameter versions of xLSTM, RWKV-4, Llama, and Mamba generalize to longer sequence lengths. Each model is trained with a context window of 2048 tokens, but during evaluation, we increase the context length up to 16384. Performance metrics are shown in Figure~\ref{fig:spaj300B_seq_extrapolation}. Among the approaches evaluated, xLSTM is notable for preserving low perplexity scores as the context length grows.

\paragraph{Validation Perplexity and Downstream Tasks.}
\label{sec:spaj_downstream_eval}
Next, across every tested model scale, we assess xLSTM, RWKV-4, Llama, and Mamba on next-token prediction accuracy using the SlimPajama validation split, as well as on a suite of downstream benchmarks aimed at evaluating models’ common sense reasoning abilities.

The third column in Table~\ref{tab:lmeval} presents the perplexity scores on the validation set for each method. Among all evaluated configurations, xLSTM[1:0] and xLSTM[7:1] consistently achieve the lowest perplexity across every model size. Performance outcomes for a range of downstream tasks are shown in the remaining columns of Table~\ref{tab:lmeval}. With few exceptions, xLSTM outperforms competing approaches at all scales, with Mamba demonstrating superior results only on the ARC benchmark in certain scenarios. Additional information is available in Appendix~\ref{sec:appExpLanguage}.

\paragraph{Performance on PALOMA Language Tasks.} Next, across all tested model scales, we evaluate the ability of xLSTM, RWKV-4, Llama, and Mamba models to predict the subsequent token in PALOMA language tasks~\citep{Magnusson:23arxivshort}. We quantify this by computing the perplexity for next token prediction over 571 distinct text domains, with sources spanning from nytimes.com to r/depression on Reddit.
Table~\ref{tab:ppleval} summarizes the token prediction perplexity, organized into broader language modeling benchmarks (first seven columns) and more specific domain benchmarks (last five columns). For these language tasks, the xLSTM[1:0] variant outperforms xLSTM[7:1].

When compared to other models, xLSTM[1:0] achieves lower perplexity than Mamba in 568 of the 571 domains (99.5\%), surpasses Llama in 486 domains (85.1\%), and outperforms RWKV-4 in 570 domains (99.8\%). Further information is provided in Appendix~\ref{sec:appExpLanguage}.

\paragraph{Scaling Laws.} In our fourth analysis, we examine the power-law scaling dynamics, a method useful for predicting model performance as size increases~\citep{Kaplan:20,Brown:20short}. The scaling patterns shown in Figure~\ref{fig:temp_sclaw300B} reveal that while all evaluated architectures demonstrate analogous scaling trends, their baselines differ. Among the models, RWKV-4 shows the poorest results, with Llama and Mamba performing incrementally better. xLSTM outperforms Mamba by a margin roughly comparable to the difference between Mamba and Llama. These scaling characteristics suggest that as model size grows, xLSTM is expected to retain its comparative advantage over Transformer and State-Space based designs.

\textbf{Generation Times and Maximal Throughput.} We evaluate the speed of text generation for various xLSTM configurations at the 1.3B parameter scale in Figure~\ref{fig:inference_time_generation}, as well as the peak throughput in Figure~\ref{fig:inference_time_decoding_throughput} (left). Our assessment includes comparable implementations of Mamba, Llama, and RWKV from HuggingFace, where the Llama model utilizes a static key-value cache. Notably, at the time of these trials, full cache compilation for the Transformer and \texttt{torch.compile} support for Mamba were unavailable. Each model undergoes testing under identical conditions with a batch size of 1 and a pre-fill length of 16 tokens—a scenario considered particularly advantageous for Transformer-based models.

As illustrated in Figure~\ref{fig:inference_time_generation}, both xLSTM and other recurrent approaches, including Mamba and RWKV-4, exhibit linear time scaling, whereas Llama displays a quadratic pattern. To examine decoding throughput, we vary the batch size and prefill parameter for the Llama model. Figure~\ref{fig:inference_time_decoding_throughput} (right) demonstrates that xLSTM sustains much higher batch sizes than Llama thanks to its constant memory footprint, resulting in superior overall throughput.

\section{Limitations}

(i) Unlike mLSTM, sLSTM’s mechanism of memory mixing restricts the use of parallelizable computations, preventing efficient parallel execution. Nonetheless, we constructed a high-speed CUDA kernel for sLSTM, which presently is less than twice as slow compared to our parallelized mLSTM kernel. (ii) The CUDA kernels currently employed for mLSTM lack proper optimization, resulting in performance that is nearly four times slower than FlashAttention or the scan utilized in Mamba. Implementation of optimized CUDA kernels in a manner similar to FlashAttention would likely yield significant speed improvements. (iii) mLSTM’s matrix memory incurs considerable computational complexity since it requires manipulating $d \times d$ matrices. However, as memory retrieval and update processes are parameter-free and compatible with conventional matrix operations, this complexity causes minimal additional wall-clock time. (iv) Proper selection of initial values for the forget gates is necessary. (v) Since matrix memory does not scale with sequence length, longer sequences may exceed memory capacity; however, up to 16k contexts this does not seem problematic, as discussed in Section~\ref{sec:spaj300B_ctx_extrapolation}. (vi) Given the high computational demands of large-scale language experiments, neither the architecture nor the hyperparameters have been thoroughly optimized, particularly for larger xLSTM configurations. A comprehensive optimization routine will be required for xLSTM models to maximize their efficacy.

\section{Conclusion}
Our investigation has partially resolved the core question: What are the limits of language modeling when LSTM architectures are scaled to billions of parameters? Current evidence suggests that such scaling allows performance to reach levels comparable to leading frameworks like Transformers and State Space Models. We proposed an extension, xLSTM, which introduces exponential gating alongside memory mixing and an updated memory structure. Compared to state-of-the-art approaches—specifically Transformers and State Space Models—xLSTM demonstrates strong results in language modeling tasks. Analysis of scaling laws reveals that increasing the size of xLSTM models positions them as formidable contenders to Transformer-based Large Language Models. Moreover, xLSTM shows promise for broad applicability, including areas such as Reinforcement Learning, Time Series Prediction, and the simulation of physical processes.

\end{document}